% Zone „Das kennst du schon“ – aus bank/geraden/zone.jsonl (zusammenbau v0.1)
\einheitenkopf[zone][Kennst du schon]{Das kennst du schon}
\setcounter{aufgabe}{0}

%% TODO Titel: Ich-kann-Satz fehlt in der Bank (Platzhalter: Kettenname) – Nr. 1
%% TODO Anweisung: ein Satz über der ersten Teilaufgabe fehlt in der Bank – Nr. 1
\begin{aufgabe}{Punkte im räumlichen Koordinatensystem lesen und Lagen aus Koordinaten deuten}
\begin{teile}
\teil Gib an, in welcher Koordinatenebene der Punkt $P(3 | 0 | 4)$ liegt.
\teil Ein Punkt liegt $2{,}5$ Einheiten unter der $x_1x_2$-Ebene, genau unter dem Punkt $(4 | -3 | 0)$. Gib seine Koordinaten an. ( \leerfeld | \leerfeld | \leerfeld )
\end{teile}
\end{aufgabe}

%% TODO Titel: Ich-kann-Satz fehlt in der Bank (Platzhalter: Kettenname) – Nr. 2
%% TODO Anweisung: ein Satz über der ersten Teilaufgabe fehlt in der Bank – Nr. 2
\begin{aufgabe}{Verbindungsvektor bilden, Vielfache erkennen, den Betrag berechnen und Vektoren normieren}
\begin{teile}
\teil Berechne den Vektor $\overrightarrow{AB}$ für $A(1 | 3 | 2)$ und $B(4 | 5 | 2)$. ( \leerfeld | \leerfeld | \leerfeld )
\teil Gib den Vektor der Länge 1 an, der in dieselbe Richtung zeigt wie $\vec u = (0 | 3 | -4)$. ( \leerfeld | \leerfeld | \leerfeld )
\end{teile}
\end{aufgabe}

%% TODO Titel: Ich-kann-Satz fehlt in der Bank (Platzhalter: Kettenname) – Nr. 3
%% TODO Anweisung: ein Satz über der ersten Teilaufgabe fehlt in der Bank – Nr. 3
\begin{aufgabe}{Skalarprodukt berechnen und Orthogonalität erkennen}
\begin{teile}
\teil Berechne das Skalarprodukt $(2 | 1 | 3) \cdot (1 | 4 | 2)$.
\end{teile}
\begin{gleichungsraster}[2]
\gl{\text{Bestimme $a$ so, dass $(4 | a | -1)$ senkrecht zu $(0{,}5 | 2 | 3)$ steht.}} & \\
\end{gleichungsraster}
\end{aufgabe}

%% TODO Titel: Ich-kann-Satz fehlt in der Bank (Platzhalter: Kettenname) – Nr. 4
%% TODO Anweisung: ein Satz über der ersten Teilaufgabe fehlt in der Bank – Nr. 4
\begin{aufgabe}{Lineare Gleichungen lösen und Widersprüche erkennen}
\begin{gleichungsraster}[2]
\gl{\text{Löse $4 + 2t = 10$.}} &
\gl{\text{Löse $2 - 4r = 8$.}} \\
\end{gleichungsraster}
\end{aufgabe}

%% TODO Titel: Ich-kann-Satz fehlt in der Bank (Platzhalter: Kettenname) – Nr. 5
%% TODO Anweisung: ein Satz über der ersten Teilaufgabe fehlt in der Bank – Nr. 5
\begin{aufgabe}{Kollinearität prüfen}
\begin{teile}
\teil Entscheide, ob $\vec a = (1 | -2 | 3)$ und $\vec b = (3 | -6 | 9)$ kollinear sind.
\teil Entscheide, ob $(0{,}5 | -1{,}5 | 2)$ und $(-2 | 6 | -8)$ kollinear sind.
\end{teile}
\end{aufgabe}

%% TODO Titel: Ich-kann-Satz fehlt in der Bank (Platzhalter: Kettenname) – Nr. 6
%% TODO Anweisung: ein Satz über der ersten Teilaufgabe fehlt in der Bank – Nr. 6
\begin{aufgabe}{Die Gerade als Gleichung in der Ebene lesen (y = m · x + n, Anstieg)}
\begin{teile}
\teil Gib Anstieg $m$ und $y$-Achsenabschnitt $n$ von $y = 3x - 2$ an. m = \leerfeld , n = \leerfeld
\teil Eine Gerade hat den Anstieg $-0{,}5$ und geht durch $(4 | 1)$. Gib ihre Gleichung an. y = \leerfeld
\end{teile}
\end{aufgabe}

%% TODO Titel: Ich-kann-Satz fehlt in der Bank (Platzhalter: Kettenname) – Nr. 7
%% TODO Anweisung: ein Satz über der ersten Teilaufgabe fehlt in der Bank – Nr. 7
\begin{aufgabe}{Prozentrechnung und Maßstab}
\begin{teile}
\teil Eine Rampe steigt auf 20 m waagerechter Länge um 1 m. Wie viel Prozent Steigung hat sie? \leerfeld \%
\teil Eine Längeneinheit entspricht 100 m. Ein Tunnel ist $12{,}5$ Längeneinheiten lang. Wie lang ist er in Kilometern? \leerfeld[km]
\end{teile}
\end{aufgabe}

%% TODO Titel: Ich-kann-Satz fehlt in der Bank (Platzhalter: Kettenname) – Nr. 8
%% TODO Anweisung: ein Satz über der ersten Teilaufgabe fehlt in der Bank – Nr. 8
\begin{aufgabe}{Lineare Gleichungen lösen und Widersprüche erkennen – Fehler finden}
\begin{teile}
\teil Finde den Fehler. Gesucht ist eine Zahl $t$, die alle drei Gleichungen erfüllt: $2 + 3t = 11$, $\quad 5 - t = 2$, $\quad -1 + 2t = 4$. \rechnung{2 + 3t &= 11 &&\mid -2 \\ t &= 3 \\ 5 - 3 &= 2 && \text{stimmt}} Antwort: $t = 3$ erfüllt alle drei Gleichungen.
\end{teile}
\end{aufgabe}

%% TODO Titel: Ich-kann-Satz fehlt in der Bank (Platzhalter: Kettenname) – Nr. 9
%% TODO Anweisung: ein Satz über der ersten Teilaufgabe fehlt in der Bank – Nr. 9
\begin{aufgabe}{Lineare Gleichungen lösen und Widersprüche erkennen – selbst rechnen}
\begin{teile}
\teil Prüfe, ob es eine Zahl $s$ gibt, die alle drei Gleichungen erfüllt: $1 + 2s = 7$, $\quad 4 - s = 1$, $\quad 3s - 2 = 8$.
\end{teile}
\end{aufgabe}
